
(en se plaçant
dans @@M0@@ grâce aux identifications 1.4).

Compte tenu de 1.6.ii), @@M1@@, et les @@M2@@ définissent un élément @@M3@@ de @@M4@@ puisque c'est un groupe séparé complet~; on posera donc @@M5@@.

Soient enfin @@M6@@, et @@M7@@~; on pose @@M8@@, @@M9@@ étant l'homomorphisme de @@M10@@ dans @@M11@@ défini en 1.2. Comme on a déjà défini @@M12@@ sur les éléments @@M13@@, il faut vérifier que les deux définitions donnent la même chose, c'est à dire que~:
@@M14@@
Or pour tout @@M15@@, on a~: @@M16@@. D'où pour tout @@M17@@~; @@M18@@. Utilisant les identifications de 1.4 et le recollement des @@M19@@ pour @@M20@@ variable, on obtient~: @@M21@@.
On a donc finalement~:
@@M22@@
Comme @@M23@@ est séparé, on en déduit l'égalité cherchée.

On a donc défini ainsi un système d'applications @@M24@@~; pour vérifier qu'elles constituent un morphisme de foncteurs (à valeurs dans @@M25@@), il faut montrer que pour tout morphisme @@M26@@, le diagramme~:
@@M27@@
